%==============================================================================
% FIG. -- robot and measurement chain (camera-ready, reviewer R2: "no image of
% the robot nor of the velocity-measurement apparatus").
% Photos are optional: drop figures/robot_photo.jpg (and, if wanted,
% figures/field_photo.jpg) into the folder and they appear with the right panel
% letters; without them the figure is the measurement-chain diagram alone.
% Photo height is capped at 3 cm so that the paper stays at 6 pages (tested).
% No number here that is not in the text (60 Hz: Sec. III-C; staircase: III-D).
%==============================================================================
\IfFileExists{figures/robot_photo.jpg}{%
  \IfFileExists{figures/field_photo.jpg}{%
    \def\setupphotos{\includegraphics[width=0.47\columnwidth,height=3.0cm,keepaspectratio]{robot_photo.jpg}\hfill
      \includegraphics[width=0.47\columnwidth,height=3.0cm,keepaspectratio]{field_photo.jpg}\\[4pt]}%
    \def\setupcap{(a)~The four-wheeled omnidirectional SSL robot used in this work.
      (b)~The field and its overhead vision cameras. (c)~Measurement chain}}{%
    \def\setupphotos{\includegraphics[width=0.60\columnwidth,height=3.0cm,keepaspectratio]{robot_photo.jpg}\\[4pt]}%
    \def\setupcap{(a)~The four-wheeled omnidirectional SSL robot used in this work.
      (b)~Measurement chain}}}{%
  \def\setupphotos{}%
  \def\setupcap{Measurement chain}}
\begin{figure}[!t]
  \centering
  \setupphotos
  \resizebox{\columnwidth}{!}{%
  \begin{tikzpicture}[
    every node/.style={font=\scriptsize},
    blk/.style={draw, rounded corners=2pt, align=center, inner sep=2.5pt, minimum height=0.72cm},
    aux/.style={blk, dashed},
    ar/.style={-{Latex[length=1.5mm]}, semithick},
    arD/.style={-{Latex[length=1.5mm]}, semithick, dashed},
    lbl/.style={font=\scriptsize\itshape, inner sep=1.5pt}
  ]
    \node[blk, text width=1.55cm] (cam)   at (0.00,0) {SSL-Vision\\ overhead cameras};
    \node[blk, text width=1.55cm] (sw)    at (2.45,0) {Team software\\ (strategy)};
    \node[blk, text width=1.25cm] (radio) at (5.15,0) {Radio\\ link};
    \node[blk, text width=1.75cm] (robot) at (7.55,0) {Robot firmware:\\ per-wheel PI loops};
    \draw[ar] (cam) -- (sw) node[lbl, midway, above] {pose};
    \draw[ar] (sw) -- (radio) node[lbl, midway, above] {$v_x, v_y, \omega$};
    \draw[ar] (radio) -- (robot);
    \draw[ar] (robot.north) -- ++(0,0.42) -| (cam.north)
      node[lbl, pos=0.25, above] {motion observed on the field};
    \node[aux, text width=2.05cm] (stair) at (2.70,-1.50)
      {Dedicated yaw collection: scripted staircase, $v_x = v_y = 0$};
    \draw[arD] (stair.north -| sw.south) -- (sw.south) node[lbl, midway, right] {replaces the strategy};
    \node[aux, text width=1.75cm] (log) at (0.00,-1.50) {Logger:\\ commands + poses};
    \draw[arD] (sw.south west) -- (log.north east);
    \node[aux, text width=2.45cm] (off) at (6.30,-1.50)
      {Offline: smoothed derivative of the pose $\rightarrow$ body velocity};
    \draw[arD] (log.south) -- ++(0,-0.30) -| (off.south);
  \end{tikzpicture}}
  \caption[Robot and measurement chain]{\setupcap: the league vision system reports the field-frame pose at about \SI{60}{Hz}; the team software closes the loop and sends body-frame velocity commands over the radio to per-wheel loops on the robot; a logger records the commands and the poses, which is all the identification consumes. Body velocities are derived offline from the vision pose (Sec.~\ref{sec:method:data}). For the dedicated yaw collection a scripted staircase of yaw commands replaces the strategy, open loop.}
  \label{fig:setup}
\end{figure}
